\documentclass[11pt,a4paper]{article}
\usepackage[margin=27mm,headheight=14pt]{geometry}
\usepackage[T1]{fontenc}
\usepackage{amsmath,amssymb,amsthm,mathtools}
\usepackage{newtxtext,newtxmath}
\usepackage{microtype,enumitem,fancyhdr,needspace}
\usepackage[hidelinks]{hyperref}
\hypersetup{pdftitle={Stochastic Online Vector Discrepancy under Causal Revelation},
  pdfsubject={A complete logarithmic-discrepancy proof against a sign-adaptive, activation-oblivious adversary},
  pdfkeywords={online discrepancy, independent activation, adaptive order, martingale, Freedman inequality}}
\pagestyle{fancy}
\fancyhf{}
\fancyhead[L]{\small Stochastic Online Vector Discrepancy}
\fancyhead[R]{\small Causal Revelation}
\fancyfoot[C]{\thepage}
\renewcommand{\headrulewidth}{0.3pt}
\setlength{\parskip}{3pt}
\setlength{\emergencystretch}{1.5em}
\setlist[enumerate]{itemsep=3pt,topsep=5pt,leftmargin=*}
\allowdisplaybreaks[1]
\numberwithin{equation}{section}
\newtheorem{theorem}{Theorem}[section]
\newtheorem{lemma}[theorem]{Lemma}
\newtheorem{proposition}[theorem]{Proposition}
\newtheorem{corollary}[theorem]{Corollary}
\theoremstyle{remark}
\newtheorem{remark}[theorem]{Remark}
\newcommand{\R}{\mathbb R}
\newcommand{\E}{\mathbb E}
\newcommand{\Pp}{\mathbb P}
\newcommand{\F}{\mathcal F}
\newcommand{\HH}{\mathcal H}
\newcommand{\ind}{\mathbf 1}
\DeclareMathOperator{\Var}{Var}
\DeclareMathOperator{\op}{op}
\title{\vspace{-1.2em}Stochastic Online Vector Discrepancy\\
  \large A complete proof under causal activation revelation}
\author{}
\date{September 29, 2026}

\begin{document}
\maketitle
\vspace{-1.8em}
\begin{abstract}
Fix known vectors $v_1,\ldots,v_n\in\R^d$ with $\|v_i\|_2\le1$ and
independent activations $A_i\sim\operatorname{Bernoulli}(p_i)$, where
$p_i\in[0,1]$. An adversary chooses the next unprocessed candidate using
past observations and signs, before inspecting that candidate's activation.
Both parties then observe the activation, and an active candidate must
immediately receive a sign. We give an online algorithm with final
discrepancy less than
$\frac{16}{3}\log\!\bigl(2(n+d)/\delta\bigr)$ with probability at least
$1-\delta$. In particular, for $d=n$ the bound is
$\frac{16}{3}\log(4n/\delta)$. The proof uses a remaining-input response
matrix, a stopped normalized martingale, an inverse-matrix telescoping
identity, and scalar Freedman concentration. An appendix proves the
finite-horizon maximal inequality used in the argument.
\end{abstract}

\section{Model and main result}
\label{sec:model}

Fix integers $n,d\ge1$. The candidate vectors and activation probabilities
are deterministic, known to the algorithm, and satisfy
\begin{equation}
 v_i\in\R^d,\qquad \|v_i\|_2\le1,\qquad p_i\in[0,1],
 \qquad i\in[n].
\end{equation}
The variables $A_1,\ldots,A_n$ are independent, with
$A_i\sim\operatorname{Bernoulli}(p_i)$. They may be sampled before the
process starts, but neither party may inspect an unrevealed activation.
Write $\mathcal A=\{i:A_i=1\}$ for the random active set.

Initially $U_0=[n]$. At step $t\in[n]$:
\begin{enumerate}
 \item The adversary selects a label $i_t\in U_{t-1}$ using the known
 data and the history of previously revealed activations and signs.
 \item Both parties observe $A_{i_t}$.
 \item If $A_{i_t}=1$, the algorithm irrevocably chooses
 $\varepsilon_{i_t}\in\{-1,+1\}$ before the next candidate is selected.
 \item Set $U_t=U_{t-1}\setminus\{i_t\}$.
\end{enumerate}
Every candidate is processed exactly once, including inactive candidates.
For an inactive candidate, set $\varepsilon_i=+1$ by convention; this has
no effect on the discrepancy. The objective is the \emph{final} norm
\begin{equation}\label{eq:discrepancy}
 \left\|\sum_{i\in\mathcal A}\varepsilon_i v_i\right\|_\infty
 =\left\|\sum_{i=1}^n A_i\varepsilon_i v_i\right\|_\infty.
\end{equation}
The adversary may react to every past sign. It may not select the next
label after inspecting its activation, and it may not inspect the
algorithm's future random choices.

\subsection{The information structure}

Let $S$ be an independent random seed for the adversary, and let
$\zeta_1,\ldots,\zeta_n$ be independent uniform random variables on $[0,1]$
used by the algorithm. The seed, the uniforms, and the activations are
mutually independent. Use the filtration
\begin{equation}\label{eq:filtration}
 \F_0=\sigma(S),\qquad
 \F_t=\sigma\!\left(S,(i_s,A_{i_s},\zeta_s)_{1\le s\le t}\right).
\end{equation}
An admissible reveal rule has $i_t\in U_{t-1}$ measurable with respect to
$\F_{t-1}$. The algorithm uses only available observations and its fresh
uniform $\zeta_t$. Including past uniforms in the analysis filtration
is harmless; the adversary need only observe the public history.

\begin{lemma}[Unrevealed activations retain their product law]
\label{lem:product}
Conditional on $\F_{t-1}$, the activations with labels in $U_{t-1}$
remain independent Bernoulli variables with their original parameters.
In particular,
\begin{equation}\label{eq:conditional-activation}
 \Pp(A_{i_t}=1\mid\F_{t-1})=p_{i_t}.
\end{equation}
\end{lemma}
\begin{proof}
Initially this follows from independence of $S$ and the activations.
Inductively, conditional on the history, the chosen label $i_t$ is fixed
before any remaining activation is inspected. Revealing $A_{i_t}$
therefore leaves the other activations with their product law. The fresh
uniform $\zeta_t$ is independent of these variables, and the sign is a
function of the existing history, $A_{i_t}$, and $\zeta_t$. Thus the new
history gives no further information about the remaining activations.
This proves the induction, including the endpoint cases $p_i=0$ and
$p_i=1$.
\end{proof}

\subsection{A theorem for finitely many test directions}

Let $\HH$ be a deterministic, nonempty finite set satisfying
\begin{equation}\label{eq:test-set}
 \{v_1,\ldots,v_n\}\subseteq\HH\subseteq
 \{h\in\R^d:\|h\|_2\le1\},\qquad L=|\HH|.
\end{equation}
Additional elements of $\HH$ specify the directions in which discrepancy
is to be controlled. All logarithms below are natural.

For $\lambda\ge1$, define
\begin{equation}\label{eq:states}
 x_t=\sum_{s=1}^t A_{i_s}\varepsilon_{i_s}v_{i_s},\qquad
 Q_t=\lambda I_d+\sum_{i\in U_t}p_i v_iv_i^\top,\qquad
 y_t=Q_t^{-1}x_t.
\end{equation}
The empty sum gives $x_0=y_0=0$. These matrices satisfy
\begin{equation}\label{eq:matrix-basic}
 Q_t\succeq\lambda I_d,\qquad
 \|Q_t^{-1}\|_{\op}\le\lambda^{-1},\qquad
 Q_t=Q_{t-1}-p_{i_t}v_{i_t}v_{i_t}^\top,\qquad
 Q_n=\lambda I_d.
\end{equation}
In particular, $Q_t$ is already $\F_{t-1}$-measurable: it depends on the
chosen label but not on its activation or sign.

\Needspace{18\baselineskip}
\begin{theorem}[Logarithmic final discrepancy]
\label{thm:main}
Fix $0<\delta<1$ and choose
\begin{equation}\label{eq:lambda}
 \lambda\ge\frac{16}{3}\log\frac{2L}{\delta}.
\end{equation}
For every admissible adaptive reveal rule, the online algorithm in
Section~\ref{sec:algorithm} satisfies
\begin{equation}\label{eq:uniform-y}
 \Pp\!\left(\max_{0\le t\le n}\max_{h\in\HH}|h^\top y_t|<1\right)
 \ge1-\delta.
\end{equation}
Consequently, with probability at least $1-\delta$,
\begin{equation}\label{eq:final-tests}
 \max_{h\in\HH}\left|\sum_{i=1}^n A_i\varepsilon_i h^\top v_i\right|
 <\lambda.
\end{equation}
The probability is over the activations and algorithmic randomness,
and also the adversary's seed if it is randomized.
\end{theorem}

Condition~\eqref{eq:lambda} implies $\lambda>1$, because $L\ge1$ and
$\delta<1$. The theorem holds separately for every admissible strategy;
it does not claim one random outcome works simultaneously for all
strategies.

\begin{corollary}[Coordinate discrepancy]
\label{cor:coordinates}
For $n$ unit-Euclidean-ball vectors in $\R^d$, there is an online signing
algorithm such that, against every admissible adversary,
\begin{equation}\label{eq:coordinate-bound}
 \Pp\!\left(
 \left\|\sum_{i\in\mathcal A}\varepsilon_i v_i\right\|_\infty
 <\frac{16}{3}\log\frac{2(n+d)}{\delta}\right)\ge1-\delta.
\end{equation}
For $d=n$, the bound is $\frac{16}{3}\log(4n/\delta)$. In this case,
for $n\ge2$ and fixed $c>0$, choosing $\delta=n^{-c}$ gives final
discrepancy $O(\log n)$ with probability at least $1-n^{-c}$.
\end{corollary}
\begin{proof}
Take $\HH=\{v_1,\ldots,v_n\}\cup\{e_1,\ldots,e_d\}$, so $L\le n+d$,
and apply Theorem~\ref{thm:main} with
$\lambda=\frac{16}{3}\log(2(n+d)/\delta)$.
\end{proof}

\section{Why the response matrix is natural}

In one dimension with every $v_i=1$, the normalized state is
\[
 y=\frac{x}{\lambda+\sum_{i\in U}p_i}.
\]
It is a slope: error per unit remaining quadratic correction capacity.
Giving the next active candidate mean sign $-y$ reduces expected error
by $p_i y$, while its revelation removes $p_i$ units of capacity.
For nonunit scalar vectors, the denominator instead contains
$\sum_{i\in U}p_i v_i^2$.

In higher dimensions, a proposed slope vector $y$ gives candidate $v_i$
the scalar sign bias $-v_i^\top y$. Its expected movement is
\begin{equation}\label{eq:response}
 -p_i(v_i^\top y)v_i=-p_i v_iv_i^\top y.
\end{equation}
Thus $v_i$ appears twice: first to measure a directional slope, then to
turn the bias into a vector movement. The linear response maps add,
giving $\sum_{i\in U}p_i v_iv_i^\top$.

If this slope were held fixed for the remaining candidates, their total
expected movement would be $-\sum_{i\in U}p_i v_iv_i^\top y$. Choosing
the slope so that the resulting residual is $\lambda y$ means
\begin{equation}\label{eq:response-balance}
 x=\lambda y+\sum_{i\in U}p_i(v_i^\top y)v_i=Qy.
\end{equation}
This explains $y=Q^{-1}x$. The term $\lambda I_d$ is a virtual residual
allowance and ensures invertibility. The function
$E_Q(x)=\frac12x^\top Q^{-1}x$ has gradient $y$, so $v_i^\top y$ is
literally its directional derivative along $v_i$.

This interpretation concerns expected movement under the stated
feedback. Individual moves need not approach zero, and an ordinary
geometric angle need not be a martingale. The argument below proves the
martingale property for the normalized state stopped at its safety boundary.

\section{The online algorithm}
\label{sec:algorithm}

Initialize $U_0=[n]$, $x_0=y_0=0$, $Q_0$ as in~\eqref{eq:states}, and a
failure flag $f_0=0$. The flag, once set, is never reset. At step $t$:
\begin{enumerate}
 \item Receive the adversary's chosen label $i=i_t\in U_{t-1}$, and compute
 \[
 r_t=v_i^\top y_{t-1},\qquad U_t=U_{t-1}\setminus\{i\},\qquad
 Q_t=Q_{t-1}-p_i v_iv_i^\top.
 \]
 \item Reveal $A_i$. If $A_i=0$ or $f_{t-1}=1$, set $\varepsilon_i=+1$.
 Otherwise use the fresh uniform $\zeta_t$ to choose
 \begin{equation}\label{eq:sign-rule}
 \varepsilon_i=
 \begin{cases}
 +1,&\zeta_t\le(1-r_t)/2,\\
 -1,&\zeta_t>(1-r_t)/2.
 \end{cases}
 \end{equation}
 Equivalently, when $f_{t-1}=0$, conditional on the history and $A_i=1$,
 $\Pp(\varepsilon_i=+1)=(1-r_t)/2$ and
 $\Pp(\varepsilon_i=-1)=(1+r_t)/2$.
 \item Update
 \[
 x_t=x_{t-1}+A_i\varepsilon_i v_i,\qquad y_t=Q_t^{-1}x_t,
 \]
 and set
 \begin{equation}\label{eq:failure-flag}
 f_t=\max\!\left\{f_{t-1},\,
 \ind_{\{\max_{h\in\HH}|h^\top y_t|\ge1\}}\right\}.
 \end{equation}
\end{enumerate}
All remaining labels are processed even after failure, using the stated
fallback sign. Hence the algorithm returns a valid signing on every
sample path.

\begin{lemma}[Validity and causality]
\label{lem:validity}
The sign rule is a valid probability distribution whenever it is used,
and the algorithm uses no unrevealed activation.
\end{lemma}
\begin{proof}
If $f_{t-1}=0$, then $|h^\top y_{t-1}|<1$ for every $h\in\HH$.
Since $v_i\in\HH$, we have $|r_t|<1$. Thus the two probabilities are
nonnegative and sum to one. The quantities $r_t$ and $Q_t$ depend only
on the known data and past history, including the current chosen label.
The sign additionally uses the revealed activation and fresh randomness.
The fallback choice is always valid.
\end{proof}

\section{Conditional noise and the stopped martingale}

\subsection{The one-step identity}

On every sample path, including after failure, define
\begin{equation}\label{eq:noise-definition}
 W_t=A_{i_t}\varepsilon_{i_t},\qquad
 \mu_t=\E[W_t\mid\F_{t-1}],\qquad \xi_t=W_t-\mu_t.
\end{equation}

\begin{lemma}[Noise moments]
\label{lem:moments}
At every step,
\begin{equation}\label{eq:noise-bounds}
 \E[\xi_t\mid\F_{t-1}]=0,\qquad
 |\xi_t|\le2,\qquad
 \E[\xi_t^2\mid\F_{t-1}]\le p_{i_t}.
\end{equation}
On the event $\{f_{t-1}=0\}$, additionally
\begin{equation}\label{eq:safe-noise}
 \mu_t=-p_{i_t}r_t,\qquad
 \xi_t=A_{i_t}\varepsilon_{i_t}+p_{i_t}r_t,\qquad
 \E[\xi_t^2\mid\F_{t-1}]=p_{i_t}-p_{i_t}^2r_t^2.
\end{equation}
\end{lemma}
\begin{proof}
Centering is immediate from~\eqref{eq:noise-definition}.
Since $W_t\in\{-1,0,1\}$, both $W_t$ and $\mu_t$ lie in $[-1,1]$,
so $|\xi_t|\le2$. Moreover, $W_t^2=A_{i_t}$, including under the fallback
rule. Lemma~\ref{lem:product} gives
\[
 \E[\xi_t^2\mid\F_{t-1}]
 =\Var(W_t\mid\F_{t-1})
 =p_{i_t}-\mu_t^2\le p_{i_t}.
\]
Before failure, the active-sign mean from~\eqref{eq:sign-rule} is $-r_t$.
Together with~\eqref{eq:conditional-activation}, this yields
$\mu_t=-p_{i_t}r_t$ and the remaining formulas. When $p_{i_t}=0$ these
identities hold directly, without conditioning on a null activation event.
\end{proof}

\begin{lemma}[Normalized update]
\label{lem:update}
On $\{f_{t-1}=0\}$,
\begin{equation}\label{eq:normalized-update}
 y_t-y_{t-1}=\xi_tQ_t^{-1}v_{i_t}.
\end{equation}
\end{lemma}
\begin{proof}
Write $p=p_{i_t}$ and $v=v_{i_t}$. Since $r_t=v^\top y_{t-1}$,
\begin{align*}
 x_t
 &=x_{t-1}+W_t v\\
 &=Q_{t-1}y_{t-1}-p(v^\top y_{t-1})v+\xi_t v\\
 &=(Q_{t-1}-pvv^\top)y_{t-1}+\xi_t v
 =Q_t y_{t-1}+\xi_t v.
\end{align*}
Multiplying by $Q_t^{-1}$ proves the claim. Because $Q_t^{-1}v$ is
$\F_{t-1}$-measurable, the conditional mean of the increment is zero
at every pre-failure state.
\end{proof}

\subsection{Stopping on the original probability space}

Let
\begin{equation}\label{eq:stopping}
 \tau=\inf\!\left\{t\in\{1,\ldots,n\}:
       \max_{h\in\HH}|h^\top y_t|\ge1\right\},
\end{equation}
with $\inf\varnothing=\infty$. Define
\begin{equation}\label{eq:stopped-states}
 \bar y_t=y_{t\wedge\tau},\qquad
 \bar Q_t=Q_{t\wedge\tau},\qquad
 B_t=\ind_{\{\tau\ge t\}}=\ind_{\{f_{t-1}=0\}}.
\end{equation}
The indicator $B_t$ is $\F_{t-1}$-measurable. In particular, the first
boundary-crossing increment is retained; only subsequent increments
are frozen.

\begin{proposition}[Stopped martingale]
\label{prop:martingale}
The process $(\bar y_t)_{t=0}^n$ is a square-integrable vector-valued
martingale with
\begin{equation}\label{eq:stopped-update}
 \bar y_t-\bar y_{t-1}=B_t\xi_tQ_t^{-1}v_{i_t}.
\end{equation}
Consequently, for every fixed $h\in\HH$,
$M_{h,t}=h^\top\bar y_t$ is a square-integrable scalar martingale
starting at zero.
\end{proposition}
\begin{proof}
On $\{B_t=1\}$, Lemma~\ref{lem:update} applies; on $\{B_t=0\}$ the
stopped increment vanishes. This proves~\eqref{eq:stopped-update}.
Its conditional mean is zero by predictability of $B_tQ_t^{-1}v_{i_t}$
and Lemma~\ref{lem:moments}. Also
$\|\bar y_t-\bar y_{t-1}\|_2\le2/\lambda$, so
$\|\bar y_t\|_2\le2n/\lambda$. This gives square integrability and
the scalar conclusion.
\end{proof}

We have not conditioned on the future success event $\{\tau>n\}$.
All conditional expectations are taken with respect to the original
filtration; the stopped process is a martingale on the full probability
space.

\section{A pathwise variance bound}

\begin{lemma}[Inverse-matrix increment]
\label{lem:inverse}
At a step with $p=p_{i_t}$ and $v=v_{i_t}$, write
$z=Q_t^{-1}v$ and $a=v^\top Q_t^{-1}v$. Then
\begin{equation}\label{eq:inverse-identity}
 Q_t^{-1}-Q_{t-1}^{-1}=\frac{p}{1+pa}zz^\top\succeq0,
\end{equation}
and
\begin{equation}\label{eq:inverse-bound}
 pzz^\top\preceq(1+\lambda^{-1})
       (Q_t^{-1}-Q_{t-1}^{-1})
 \preceq2(Q_t^{-1}-Q_{t-1}^{-1}).
\end{equation}
\end{lemma}
\begin{proof}
We have $Q_{t-1}=Q_t+pvv^\top$. Multiplication verifies
\[
 (Q_t+pvv^\top)^{-1}
 =Q_t^{-1}-\frac{p}{1+pa}zz^\top,
\]
which is also the rank-one inverse formula. This proves
\eqref{eq:inverse-identity}, including $p=0$. Since
$0\le a\le\|v\|_2^2/\lambda\le1/\lambda$ and $p\le1$, multiplying
the identity by $1+pa\le1+1/\lambda\le2$ gives
\eqref{eq:inverse-bound}.
\end{proof}

For fixed $h\in\HH$, let
\begin{equation}\label{eq:scalar-variance}
 D_{h,t}=M_{h,t}-M_{h,t-1},\qquad
 V_{h,t}=\sum_{s=1}^t\E[D_{h,s}^2\mid\F_{s-1}].
\end{equation}

\begin{proposition}[Deterministic increment and variance budgets]
\label{prop:variance}
For every $h\in\HH$, almost surely,
\begin{equation}\label{eq:variance-budget}
 |D_{h,t}|\le\frac2\lambda\quad(1\le t\le n),\qquad
 V_{h,n}\le\frac{1+1/\lambda}{\lambda}\le\frac2\lambda.
\end{equation}
\end{proposition}
\begin{proof}
Equation~\eqref{eq:stopped-update} gives
$D_{h,t}=B_t\xi_t h^\top Q_t^{-1}v_{i_t}$, so its absolute value is at
most $2\|h\|_2\|v_{i_t}\|_2/\lambda\le2/\lambda$.
Using $B_t^2=B_t$, predictability, and the conditional noise bound,
\begin{align*}
 \E[D_{h,t}^2\mid\F_{t-1}]
 &\le B_t p_{i_t}(h^\top Q_t^{-1}v_{i_t})^2\\
 &\le(1+1/\lambda)B_t
       h^\top(Q_t^{-1}-Q_{t-1}^{-1})h.
\end{align*}
The second inequality is Lemma~\ref{lem:inverse}. Pathwise,
\[
 B_t(Q_t^{-1}-Q_{t-1}^{-1})
 =\bar Q_t^{-1}-\bar Q_{t-1}^{-1}.
\]
Summing and using $\bar Q_n\succeq\lambda I_d$ yields
\begin{align*}
 V_{h,n}
 &\le(1+1/\lambda)h^\top(\bar Q_n^{-1}-Q_0^{-1})h\\
 &\le(1+1/\lambda)h^\top\bar Q_n^{-1}h
 \le\frac{1+1/\lambda}{\lambda}\|h\|_2^2
 \le\frac2\lambda.
\end{align*}
This is a pathwise bound, valid even when failure occurs.
\end{proof}

\Needspace{18\baselineskip}
\section{Concentration and completion of the proof}

We use the following finite-horizon form of Freedman's inequality
\cite{freedman1975}. Appendix~\ref{app:freedman} proves it directly.

\begin{lemma}[Two-sided maximal Freedman inequality]
\label{lem:freedman}
Let $(M_t,\F_t)_{t=0}^n$ be a square-integrable scalar martingale with
$M_0=0$, and set $D_t=M_t-M_{t-1}$. Suppose deterministic constants
$R>0$ and $\nu\ge0$ satisfy, almost surely,
\[
 |D_t|\le R\quad(1\le t\le n),\qquad
 \sum_{t=1}^n\E[D_t^2\mid\F_{t-1}]\le\nu.
\]
Then, for every $b>0$,
\begin{equation}\label{eq:freedman}
 \Pp\!\left(\max_{0\le t\le n}|M_t|\ge b\right)
 \le2\exp\!\left(-\frac{b^2}{2(\nu+Rb/3)}\right).
\end{equation}
\end{lemma}

\begin{proof}[Proof of Theorem~\ref{thm:main}]
For a fixed $h\in\HH$, Propositions~\ref{prop:martingale}
and~\ref{prop:variance} allow Lemma~\ref{lem:freedman} with
$R=2/\lambda$, $\nu=2/\lambda$, and $b=1$. Hence
\begin{equation}\label{eq:one-direction-tail}
 \Pp\!\left(\max_{0\le t\le n}|h^\top\bar y_t|\ge1\right)
 \le2\exp\!\left(-\frac{1}{2(2/\lambda+2/(3\lambda))}\right)
 =2e^{-3\lambda/16}.
\end{equation}
If $\tau\le n$, some $h\in\HH$ satisfies $|h^\top y_\tau|\ge1$, and
$\bar y_\tau=y_\tau$. Therefore
\begin{equation}\label{eq:failure-probability}
 \Pp(\tau\le n)
 \le\sum_{h\in\HH}
      \Pp\!\left(\max_{0\le t\le n}|h^\top\bar y_t|\ge1\right)
 \le2L e^{-3\lambda/16}\le\delta.
\end{equation}
The final inequality is~\eqref{eq:lambda}. On $\{\tau>n\}$ the safety
boundary is never reached, proving~\eqref{eq:uniform-y}.

Every label is processed once, so $U_n=\varnothing$ and
$Q_n=\lambda I_d$ on every sample path. Consequently, on the same
success event,
\[
 \max_{h\in\HH}|h^\top x_n|
 =\lambda\max_{h\in\HH}|h^\top y_n|<\lambda.
\]
Substituting $x_n=\sum_i A_i\varepsilon_i v_i$ proves
\eqref{eq:final-tests}.
\end{proof}

\begin{remark}[No union bound over time]
The maximal inequality already controls every time for a fixed test
direction. The only union bound in~\eqref{eq:failure-probability} is
over the $L$ fixed directions.
\end{remark}

\begin{remark}[Scope of the guarantee]
The theorem bounds $y_t$ in the test directions at every time and bounds
the original discrepancy $x_n$ at completion. It does not assert the
same bound for every prefix $x_t$. Processing inactive candidates is
essential to this argument: it preserves the stated information model and makes
$Q_n=\lambda I_d$. If the adversary can inspect unrevealed activations,
Lemma~\ref{lem:product} need not hold. The argument also requires
independent activation variables; candidates sharing a single random
activation cannot automatically be treated as independent.
\end{remark}

\begin{remark}[Scaling and zero probabilities]
If instead $\|v_i\|_2\le\rho$ for a fixed $\rho>0$, apply the algorithm
to $v_i/\rho$ and multiply the discrepancy bound by $\rho$.
Candidates with $p_i=0$ may be retained or deleted in advance.
The theorem requires neither a positive lower bound on the probabilities
nor an upper bound strictly below one. No optimality claim is made for
the logarithmic upper bound.
\end{remark}

\clearpage
\appendix
\section{A proof of the finite-horizon Freedman inequality}
\label{app:freedman}

We prove Lemma~\ref{lem:freedman}. Write
\[
 q_t=\E[D_t^2\mid\F_{t-1}],\qquad V_t=\sum_{s=1}^t q_s.
\]
Then $V_n\le\nu$ almost surely, $q_t$ is predictable, and
$\E[D_t\mid\F_{t-1}]=0$.

\subsection{A conditional exponential-moment bound}

For every integer $k\ge2$, $k!\ge2\cdot3^{k-2}$. Thus, for $|u|<3$,
\begin{equation}\label{eq:taylor}
 e^u
 =1+u+\sum_{k=2}^{\infty}\frac{u^k}{k!}
 \le1+u+\frac{u^2}{2}\sum_{j=0}^{\infty}(|u|/3)^j
 =1+u+\frac{u^2}{2(1-|u|/3)}.
\end{equation}
Fix $0<\theta<3/R$. Since $|D_t|\le R$, taking conditional
expectations in~\eqref{eq:taylor} gives
\begin{equation}\label{eq:mgf}
 \E[e^{\theta D_t}\mid\F_{t-1}]
 \le1+\frac{\theta^2 q_t}{2(1-\theta R/3)}
 \le\exp\!\bigl(g(\theta)q_t\bigr),\qquad
 g(\theta)=\frac{\theta^2}{2(1-\theta R/3)}.
\end{equation}
The last step uses $1+a\le e^a$.

\subsection{A nonnegative supermartingale and first crossing}

Set
\[
 Z_t=\exp\!\bigl(\theta M_t-g(\theta)V_t\bigr),\qquad Z_0=1.
\]
Predictability of $q_t$ and~\eqref{eq:mgf} imply
\[
 \E[Z_t\mid\F_{t-1}]
 =Z_{t-1}e^{-g(\theta)q_t}
       \E[e^{\theta D_t}\mid\F_{t-1}]
 \le Z_{t-1}.
\]
Thus $(Z_t)$ is a nonnegative supermartingale. Integrability follows
from $Z_t\le e^{\theta nR}$, since $M_t\le nR$ and $V_t\ge0$.

Let $E_b=\{\max_{0\le t\le n}M_t\ge b\}$ and
\[
 \kappa=\min\!\left\{\inf\{t\in\{0,\ldots,n\}:M_t\ge b\},\,n\right\},
\]
with $\inf\varnothing=\infty$. This is a stopping time bounded by $n$.
The pathwise telescoping identity
\[
 Z_\kappa=1+\sum_{t=1}^n
       \ind_{\{\kappa\ge t\}}(Z_t-Z_{t-1})
\]
and predictability of $\ind_{\{\kappa\ge t\}}$ show
$\E[Z_\kappa]\le1$. This uses only a finite sum.
On $E_b$, we have $M_\kappa\ge b$ and $V_\kappa\le\nu$, so
\[
 Z_\kappa\ge\exp\!\bigl(\theta b-g(\theta)\nu\bigr).
\]
Nonnegativity then gives
\begin{equation}\label{eq:crossing-bound}
 \Pp(E_b)\le\exp\!\bigl(-\theta b+g(\theta)\nu\bigr).
\end{equation}

\subsection{The parameter choice and both tails}

If $\nu>0$, choose
\[
 \theta=\frac{b}{\nu+Rb/3}.
\]
This lies in $(0,3/R)$, and
$1-\theta R/3=\nu/(\nu+Rb/3)$. Consequently,
\[
 -\theta b+g(\theta)\nu
 =-\frac{b^2}{\nu+Rb/3}
   +\frac{b^2}{2(\nu+Rb/3)}
 =-\frac{b^2}{2(\nu+Rb/3)}.
\]
Substitution in~\eqref{eq:crossing-bound} proves the upper-tail
maximal inequality. Apply the same argument to $-M_t$, which has the
same increment and variance bounds, and take a union bound over the
two tails to obtain~\eqref{eq:freedman}.

If $\nu=0$, nonnegativity of the $q_t$ implies $q_t=0$ almost surely
for every $t$. Hence $\E[D_t^2]=0$, every increment vanishes almost
surely, and the inequality follows immediately. This completes the proof.

\begin{thebibliography}{9}
\bibitem{freedman1975}
David A. Freedman.
\newblock On tail probabilities for martingales.
\newblock \emph{The Annals of Probability}, 3(1):100--118, 1975.
\newblock \href{https://doi.org/10.1214/aop/1176996452}{doi:10.1214/aop/1176996452}.
\end{thebibliography}
\end{document}
